\documentclass[11pt]{article}
\usepackage[margin=1in]{geometry}
\usepackage{amsmath,amssymb,amsthm,hyperref}
\newtheorem{theorem}{Theorem}
\newcommand{\fl}[1]{\left\lfloor #1\right\rfloor}
\title{An explicit estimate for Garcia's paired floor sum}
\author{Principia Math In Progress}
\date{29 September 2026}
\begin{document}
\maketitle

\paragraph{Task 1 (original problem).}
MathDB \href{https://mathdb.com/p/357173}{357173} records Conjecture 6.6
of Garcia \cite[p.~30]{Garcia}. For positive integers \(n\) with
\(p=4n-1\) prime, put
\[
 M=\fl{\frac{n^2}{p}},\qquad R_m=\sqrt{mp},\qquad
 Q_m=\frac{1+\sqrt{4mp+3p-4}}2,
\]
and
\[
 A_n=\sum_{m=1}^M\fl{R_m}+\sum_{m=0}^{M-1}\fl{Q_m},
 \qquad D_n=Mp+2n.
\]
Prove that \(A_n/D_n\to1/3\). The source also proposes
\(\fl{D_n/3}\) as a \emph{good estimate} without specifying an error.
The following theorem proves the limit and gives a precise relative-error
interpretation of that estimate.

\begin{theorem}
For every integer \(n\ge4\), with the definitions above and without any
primality assumption,
\[
 \left|A_n-\frac{D_n}{3}\right|\le4n,\qquad
 \left|A_n-\fl{D_n/3}\right|\le4n+1,\qquad
 \left|\frac{A_n}{D_n}-\frac13\right|\le\frac8n.
\]
In particular the conjectured limit holds along the prime subsequence.
\end{theorem}

\begin{proof}
First,
\[
 \frac{n^2}{4n-1}=\frac n4+\frac{n}{4(4n-1)},\qquad
 0<\frac{n}{4(4n-1)}<\frac14.
\]
The fractional part of \(n/4\) is one of \(0,1/4,1/2,3/4\), so
\(M=\fl{n/4}\). In particular \(M\ge1\).

Write
\[
 B=\sum_{m=1}^M\fl{\sqrt{mp}},\qquad
 T=\sqrt p\sum_{m=1}^M\sqrt m,\qquad
 V=\frac23\sqrt p\,M^{3/2}.
\]
Monotonicity of \(\sqrt x\) gives
\[
 V=\sqrt p\int_0^M\sqrt x\,dx\ \le T\le V+\sqrt{pM}.
\]
For the upper bound, apply the integral comparison to the terms with
\(m<M\), and then add \(\sqrt M\). Rounding each summand gives
\[
 T-M\le B\le T.
\]
Since \(p\ge15\), for every \(0\le m<M\),
\[
 R_m\le Q_m=\frac12+\sqrt{mp+\frac{3p}{4}-1}
          \le R_{m+1}+1.
\]
Consequently, if \(C=\sum_{m=0}^{M-1}\fl{Q_m}\), then
\[
 B-\fl{R_M}\le C\le B+M.
\]
Combining these estimates and writing \(W=2V\) proves the useful sandwich
\begin{equation}\label{sandwich}
 W-2M-\sqrt{pM}\le A_n\le W+2\sqrt{pM}+M.
\end{equation}
This argument retains every endpoint, including integral square roots.

For completeness we turn the sandwich into explicit bounds. Write
\(r=n-4M\in\{0,1,2,3\}\). Then
\[
 W=\frac{n^2}{3}
       \sqrt{1-\frac1{4n}}\left(1-\frac rn\right)^{3/2}.
\]
For \(0\le x,y\le1\),
\[
 1-\sqrt{1-x}\le x,\qquad
 1-(1-y)^{3/2}\le\frac32y .
\]
The second inequality follows by integrating the derivative
\(\frac32\sqrt{1-y}\le\frac32\). Since \(1-ab\le(1-a)+(1-b)\)
for \(a,b\in[0,1]\), we get
\[
 0\le\frac{n^2}{3}-W
 \le\frac n{12}+\frac{rn}{2}\le\frac{19n}{12}.
\]
Also \(M\le n/4\) and \(\sqrt{pM}\le n\). Thus \eqref{sandwich} yields
\[
 -\frac{37n}{12}\le A_n-\frac{n^2}{3}\le\frac{9n}{4},
 \qquad
 \left|A_n-\frac{n^2}{3}\right|\le\frac{37n}{12}.
\]
Direct expansion gives
\[
 D_n=n^2+\frac{(7-4r)n+r}{4},\qquad
 \left|\frac{D_n}{3}-\frac{n^2}{3}\right|\le\frac{2n}{3}.
\]
The last inequality holds for each of the four possible \(r\).
It follows that
\[
 \left|A_n-\frac{D_n}{3}\right|
 \le\frac{45n}{12}\le4n.
\]
Rounding \(D_n/3\) changes it by less than one, proving the second claim.
Finally,
\[
 D_n\ge n^2-\frac{5n}{4}\ge\frac{n^2}{2}\qquad(n\ge4).
\]
Dividing the first bound by \(D_n\) proves the third claim and the limit.
\end{proof}

\paragraph{Scope and verification.}
The proof uses only the definitions of the two roots, an integral
comparison, and bounded rounding errors. It does not assume any of the
source's class-number or jump-count identities. The source's qualitative
phrase \emph{good estimate} is made precise here as an \(O(n)\) absolute
error on a quantity of order \(n^2\); no bounded absolute error was
specified or is asserted.
The accompanying integer-arithmetic verifier checks the original two
sums, their root inequalities, and the explicit bounds for every
\(4\le n\le4000\), together with larger examples. These finite checks
support the proof but do not establish the infinite limit.

\paragraph{Task 2 (proposed follow-up).}
Determine the optimal constant in a uniform bound
\(\lvert A_n-D_n/3\rvert\le c n+O(1)\).

\paragraph{Task 3 (proposed follow-up).}
For each residue class of \(n\) modulo four, determine whether
\((A_n-D_n/3)/n\) converges, and compare unrestricted and prime subsequences.

\paragraph{Task 4 (proposed follow-up).}
Obtain smaller remainder terms after subtracting any second-order main
term, including the effect of square divisors of \(4n-1\).

\paragraph{Task 5 (proposed follow-up).}
Determine the corresponding limit and explicit error when
\(p/n\to c>0\), \(M=\fl{n^2/p}\), and \(Q_m\) is replaced by
\(\alpha+\sqrt{mp+\beta p+\gamma}\), for fixed real parameters whenever
all summands are defined.

\begin{thebibliography}{1}
\bibitem{Garcia}
J. Garcia, \emph{Class Number of the Imaginary Quadratic Field and
Quadratic Residues Identities}, arXiv:2301.02951v1 (2023), 31 pages,
\url{https://arxiv.org/abs/2301.02951}.
Definitions: Notation 2.2; target: Conjecture 6.6.
\end{thebibliography}
\end{document}
